\documentclass[10pt]{article}
\usepackage{hyperref}
\usepackage{enumitem}
\usepackage{cmap}
\usepackage[utf8]{inputenc}
\usepackage[T1]{fontenc}
\usepackage[english]{babel}
\usepackage[left=1.05cm,top=0.92cm,right=1.05cm,bottom=0.78cm]{geometry}
\IfFileExists{glyphtounicode.tex}{\input{glyphtounicode}\pdfgentounicode=1}{}
\hypersetup{
    hidelinks,
    pdfauthor={Kevin Bravo},
    pdftitle={Kevin Bravo - Founding Product Engineer (Mobile, AI-Native)},
    pdfsubject={CV tailored for the TAYA founding engineer role: consumer mobile product engineering, 0-to-1 building, AI-native workflows, and product ownership},
    pdfkeywords={Kevin Bravo, Founding Engineer, Product Engineer, Mobile, iOS, Swift, SwiftUI, React Native, Expo, Flutter, Dart, TypeScript, Supabase, AI-native, Claude Code, Codex, 0-to-1, consumer apps, privacy}
}
\setlength{\parindent}{0pt}
\hyphenpenalty=10000
\exhyphenpenalty=10000
\emergencystretch=2em
\pagestyle{empty}
\setlist[itemize]{leftmargin=*, noitemsep, topsep=1pt, partopsep=0pt, parsep=1pt, label={\textbf{-}}}

\newcommand{\sectiontitle}[1]{%
    \vspace{5pt}
    \begin{center}
        \textbf{#1}
    \end{center}
    \vspace{-6pt}
}

\begin{document}
\raggedright

\begin{center}
    {\LARGE \textbf{Kevin Bravo}}\\
    Founding Product Engineer / Full-Stack \& Mobile \\
    Consumer Product, 0$\rightarrow$1 Building \& AI-Native Development \\
    \hrulefill
\end{center}

\begin{center}
    Email: \href{mailto:0bravokevin@gmail.com}{0bravokevin@gmail.com}
    | Portfolio: \href{https://kevinbravo.com}{https://kevinbravo.com}
    | GitHub: \href{https://github.com/0bkevin}{https://github.com/0bkevin}
    | LinkedIn: \href{https://www.linkedin.com/in/0bkevin/}{https://www.linkedin.com/in/0bkevin}
    | Remote
\end{center}

\sectiontitle{Professional Summary}
Founding-engineer-style product engineer with 5 years of experience taking products from a rough idea to shipped software and owning them after they ship. I build across the full stack and on mobile, with cross-platform consumer apps in React Native (Expo) and Flutter, and I am actively building with Swift and SwiftUI to go deeper on native iOS. I work AI-native, using Claude Code and Codex every day, and I care a lot about how a product looks, feels, and reads. I do my best work with a founder, no ticket queue, and a clear problem where the product still needs to be figured out.

\sectiontitle{Technical Profile}
\textbf{Mobile product:} React Native, Expo, Flutter, Dart, cross-platform consumer apps, offline and local-first patterns, BLE-adjacent device sync (companion/relay protocols), background services, app packaging. Swift and SwiftUI: familiar and actively learning. \\
\textbf{Full-stack:} TypeScript, JavaScript, React, Next.js, Astro, Node.js, Python, Django, REST APIs, PostgreSQL, Prisma, SQLite, Supabase-compatible schema design, managed backends. \\
\textbf{AI-native development:} Claude Code, Codex, Cursor, OpenCode, LLM workflows, AI agents, agent tooling and control planes; building daily with coding agents while keeping engineering judgment in the loop. \\
\textbf{Product and craft:} product taste, interaction detail, privacy and data minimization, onboarding and first-run flows, 0$\rightarrow$1 delivery under ambiguity, release ownership, and production support.

\sectiontitle{Building Right Now}
\textbf{Brio} (Expo React Native, Go, Postgres): A mobile control plane for a personal AI agent. I designed it as a small distributed system: an Expo React Native app, a Go companion service, a Go relay for remote connections, and a shared JSON protocol. It handles device pairing, an ownership model for persistent agents across devices, background service installation on macOS, Linux, and Windows, and health checks. \\
\textbf{Zrode} (TypeScript): My own version of an open-source web and desktop GUI for coding agents (Codex, Claude, OpenCode), where I am shaping the workspace and agent experience the way I want to work. \\
\textbf{Hermes:} A personal AI agent running 24/7 on a VPS for tasks, research, and daily workflows.

\sectiontitle{Selected Experience}

\textbf{Proyecto Educa} \hfill \textit{Remote} \\
\textbf{Chief Technology Officer} \hfill \textit{Dec 2025 - Present}
\begin{itemize}
    \item \textbf{Own product and technical direction} end to end for an education platform, from architecture to shipped software, and I both set the direction and write the code.
    \item \textbf{Leading the design and build of an internal streaming platform} so students can access courses directly through the product.
\end{itemize}

\vspace{4pt}
\textbf{Freelance} \hfill \textit{Remote} \\
\textbf{Software Engineer (Mobile \& Full-Stack)} \hfill \textit{Sep 2024 - Dec 2024}
\begin{itemize}
    \item \textbf{Built Knowfi, a cross-platform consumer mobile app} (Flutter, Dart) for personal finance, with state-driven flows and a shared codebase for iOS and Android, taken from empty project to working product.
    \item \textbf{Shipped consumer web products end to end} (Astro), including jenndolinski.com and rodolfomedinadelrio.com, owning UI, performance, and deployment. One client, Proyecto Educa, later brought me on as CTO.
\end{itemize}

\vspace{4pt}
\textbf{Free2Z} \hfill \textit{Remote} \\
\textbf{Principal Software Engineer} \hfill \textit{Apr 2025 - Apr 2026}
\begin{itemize}
    \item \textbf{Led the migration of a live consumer product from React to SvelteKit} with real users on it, rebuilding frontend architecture and SSR without breaking a running product.
    \item \textbf{Integrated Dyte live video} and shipped a CMS suite with a Markdown editor and blog engine, working across product, implementation, and support.
\end{itemize}

\vspace{4pt}
\textbf{U.S. Department of State (via AVAA)} \hfill \textit{Remote} \\
\textbf{Software Engineer | Full-Stack Developer} \hfill \textit{Nov 2024 - Aug 2025}
\begin{itemize}
    \item \textbf{Sole engineer on Billingua Talent}, a multi-tenant platform with role-based access for four user types, built end to end from data model to interfaces and tested in pilots with 50 users and 5 companies.
    \item \textbf{Owned the whole arc}: requirements, architecture, product flows, and every release, with no project managers in between.
\end{itemize}

\vspace{4pt}
\textbf{Asociacion Venezolano Americana de Amistad (AVAA)} \hfill \textit{Remote} \\
\textbf{Software Developer} \hfill \textit{Oct 2022 - Jul 2024}
\begin{itemize}
    \item \textbf{Built and operate SEP} (Next.js, PostgreSQL, Prisma, Auth.js on Azure), a multi-tenant platform serving around 300 scholars and staff across 3 national chapters, and stayed its technical owner in production.
\end{itemize}

\sectiontitle{Why This Role}
I want to own the software at the center of a product people open every day, work directly with the founder with no layers, and measure success by whether people keep coming back. I build consumer mobile apps, I am AI-native by default, and I care about the details nobody asked about. Native iOS in Swift is where I am investing my learning now, and I would be joining to go all in on it.

\sectiontitle{Education}
\textbf{IESA} — Master of Public Administration \hfill \textit{In Progress} \\
\textbf{Universidad Central de Venezuela} — Occupational Therapy, Cum Laude \hfill \textit{2023}

\sectiontitle{Additional}
\textbf{Mobile/Product:} React Native, Expo, Flutter, Dart, Swift/SwiftUI (learning), consumer apps, privacy-minded design. \textbf{Full-stack:} TypeScript, React, Next.js, Node.js, Python, PostgreSQL, Prisma, Supabase. \textbf{AI:} Claude Code, Codex, AI agents, LLM workflows. \textbf{Languages:} Spanish (Native), English (C1).

\end{document}
